% ======================================================================
\chapter[Операторы. Собственные функции и собственные значения]%
        {Квантово-механические операторы. Собственные функции и значения}
\label{t:05}
% ======================================================================

\begin{qbox}
\textbf{Билет 5.} Понятие о квантово-механических операторах. Операторы координаты,
импульса, потенциальной и кинетической энергии системы, гамильтониан. Собственные функции и
собственные значения операторов.
\end{qbox}

\section{Содержание билета}

\textbf{Подтемы:}
\begin{enumerate}
 \item оператор, линейность, эрмитовость; связь операторов с наблюдаемыми величинами;
 \item операторы координаты, импульса, кинетической и потенциальной энергии, гамильтониан (для одной частицы и системы);
 \item задача на собственные значения; спектр (дискретный и непрерывный), вырождение;
 \item свойства собственных функций эрмитовых операторов: действительность собственных значений,
       ортогональность, полнота; вероятности результатов измерений и средние;
 \item коммутаторы, одновременная измеримость, соотношение неопределённостей.
\end{enumerate}

\textbf{Необходимые понятия:} волновая функция и интерпретация Борна (\tref{t:04});
скалярное произведение функций; интегрирование по частям; ряд и интеграл Фурье.

\section{Определения и физические принципы}

\begin{definition}[оператор]
\emph{Оператор} $\op A$ — правило, сопоставляющее функции $\psi$ другую функцию $\varphi=\op A\psi$.
Оператор \emph{линеен}, если $\op A(c_1\psi_1+c_2\psi_2)=c_1\op A\psi_1+c_2\op A\psi_2$.
Примеры: умножение на функцию $f(x)$, дифференцирование $\dd/\dd x$.
\end{definition}

\begin{definition}[эрмитов оператор]
Скалярное произведение $\langle\varphi|\psi\rangle=\int\varphi^*\psi\,\dd V$. Оператор $\op A$
\emph{эрмитов} (самосопряжён), если для любых допустимых $\varphi,\psi$
\begin{equation}\label{eq:5:herm}
  \int\varphi^*\,(\op A\psi)\,\dd V=\int(\op A\varphi)^*\,\psi\,\dd V .
\end{equation}
\end{definition}

\begin{definition}[собственные функции и значения]
Если $\op A\psi_a=a\,\psi_a$, где $a$ — число, а $\psi_a$ удовлетворяет стандартным условиям,
то $\psi_a$ — \emph{собственная функция}, $a$ — \emph{собственное значение}. Множество собственных
значений — \emph{спектр} оператора (дискретный, непрерывный или смешанный). Если одному $a$
отвечают $g>1$ линейно независимых собственных функций, значение $a$ \emph{вырождено} с
кратностью $g$.
\end{definition}

\begin{principle}{Постулаты, связывающие операторы с измерениями}
\begin{enumerate}[leftmargin=1.4em]
 \item Каждой физической величине $A$ сопоставляется линейный эрмитов оператор $\op A$.
 \item При измерении $A$ могут получиться только собственные значения $a_n$ оператора $\op A$.
 \item Если $\psi=\sum_nc_n\psi_n$ — разложение состояния по нормированным собственным функциям,
       то вероятность получить $a_n$ равна $|c_n|^2$, где $c_n=\int\psi_n^*\psi\,\dd V$.
 \item Среднее значение величины в состоянии $\psi$:
       \begin{equation}\label{eq:5:mean}
         \avg A=\int\psi^*\op A\psi\,\dd V .
       \end{equation}
 \item После измерения, давшего $a_n$, система оказывается в состоянии $\psi_n$ (для невырожденного $a_n$).
\end{enumerate}
\end{principle}

\emph{Физический смысл.} В состоянии, являющемся собственной функцией $\op A$, величина $A$ имеет
точное значение: разброс результатов $\Delta A=0$. В общем состоянии $A$ не имеет определённого
значения, и квантовая механика даёт лишь распределение вероятностей. Эрмитовость нужна для
того, чтобы результаты измерений (собственные значения) и средние были действительными.

\section{Основная теория}

\subsection{Почему величинам сопоставляются операторы}

В классической механике состояние — точка фазового пространства $(x,p)$, и любая величина —
функция $A(x,p)$. В квантовой механике состояние — волновая функция, и вопрос «каково значение
$p$ в состоянии $\psi$?» в общем случае не имеет ответа (соотношение неопределённостей). Можно
спрашивать лишь о распределении результатов. Оказывается, всю информацию о распределении
величины удобно закодировать оператором: его собственные значения — возможные результаты,
собственные функции — состояния с определённым значением, а формула~\eqref{eq:5:mean} даёт среднее.

\subsection{Основные операторы}

\paragraph{Координата.} Поскольку $|\psi(\vect r)|^2$ — плотность вероятности координаты,
$\avg{\vect r}=\int\psi^*\vect r\psi\,\dd V$, т.\,е. $\op{\vect r}=\vect r$ — оператор умножения на
координату. Аналогично $\op U=U(\vect r)$ — оператор потенциальной энергии (умножение на функцию).

\paragraph{Импульс.} Плоская волна $\ee^{\ii px/\hbar}$ — состояние с определённым импульсом $p$;
она должна быть собственной функцией $\op p_x$ с собственным значением $p$. Этому требованию
удовлетворяет дифференциальный оператор
\begin{equation}\label{eq:5:p}
  \op p_x=-\ii\hbar\pdiff{}{x},\qquad \op{\vect p}=-\ii\hbar\nabla:\qquad
  -\ii\hbar\pdiff{}{x}\ee^{\ii px/\hbar}=p\,\ee^{\ii px/\hbar}.
\end{equation}
Строгое обоснование: среднее импульса, вычисленное через распределение $|a(p)|^2$ (\tref{t:04}),
совпадает с $\int\psi^*(-\ii\hbar\,\dd/\dd x)\psi\,\dd x$ (вывод~\ref{der:5:pmean}).

\paragraph{Кинетическая энергия и гамильтониан.} Операторы функций от импульса строятся
подстановкой $\vect p\to\op{\vect p}$:
\begin{equation}\label{eq:5:T}
  \op T=\frac{\op{\vect p}^{\,2}}{2m}=-\frac{\hbar^2}{2m}\nabla^2=
  -\frac{\hbar^2}{2m}\Bigl(\pdiff{^2}{x^2}+\pdiff{^2}{y^2}+\pdiff{^2}{z^2}\Bigr).
\end{equation}
Гамильтониан — оператор полной энергии:
\begin{equation}\label{eq:5:H}
  \op H=\op T+\op U=-\frac{\hbar^2}{2m}\nabla^2+U(\vect r).
\end{equation}
Для системы $N$ частиц с массами $m_i$, взаимодействующих между собой и находящихся во внешнем поле:
\begin{equation}\label{eq:5:HN}
  \op H=\sum_{i=1}^{N}\Bigl(-\frac{\hbar^2}{2m_i}\nabla_i^2\Bigr)+U(\vect r_1,\dots,\vect r_N),
\end{equation}
где $\nabla_i$ действует на координаты $i$-й частицы. Например, для атома гелия (ядро
неподвижно) $U=-2e^2/r_1-2e^2/r_2+e^2/|\vect r_1-\vect r_2|$.

\begin{remark}[Правило соответствия]
Переход $A(\vect r,\vect p)\to\op A=A(\op{\vect r},\op{\vect p})$ однозначен, если в $A$ нет произведений
некоммутирующих множителей вида $xp_x$. В общем случае требуется симметризация, например
$xp_x\to\frac12(\op x\op p_x+\op p_x\op x)$, чтобы оператор был эрмитовым.
\end{remark}

\subsection{Свойства собственных функций эрмитовых операторов}

\begin{enumerate}
 \item \textbf{Собственные значения действительны} (вывод~\ref{der:5:real}).
 \item \textbf{Собственные функции, отвечающие разным собственным значениям, ортогональны}
       (вывод~\ref{der:5:orth}); в случае вырождения их можно выбрать ортогональными.
 \item \textbf{Полнота:} любую допустимую функцию можно разложить по собственным функциям
       эрмитова оператора (утверждение без доказательства):
       \begin{equation}\label{eq:5:expand}
         \psi=\sum_nc_n\psi_n+\int c(a)\psi_a\,\dd a,\qquad c_n=\int\psi_n^*\psi\,\dd V .
       \end{equation}
       Для нормированного $\psi$ $\sum|c_n|^2+\int|c(a)|^2\dd a=1$ — сумма вероятностей всех исходов.
\end{enumerate}
Из разложения следует, что среднее~\eqref{eq:5:mean} равно $\sum_na_n|c_n|^2$ — обычному
статистическому среднему (вывод~\ref{der:5:avg}).

\subsection{Примеры задач на собственные значения}

\textbf{Импульс} ($-\ii\hbar\psi'=p\psi$): $\psi_p=C\ee^{\ii px/\hbar}$; ограниченность при
$x\to\pm\infty$ требует действительного $p$; спектр непрерывный, $-\infty<p<\infty$; функции
нормируются на $\delta$-функцию (\tref{t:04}).

\textbf{Координата} ($x\psi=x_0\psi$): $\psi_{x_0}=\delta(x-x_0)$ — состояние, локализованное в точке;
спектр непрерывный.

\textbf{Проекция момента импульса} $\op L_z=-\ii\hbar\,\partial/\partial\varphi$: условие однозначности
$\psi(\varphi+2\pi)=\psi(\varphi)$ делает спектр \emph{дискретным}: $L_z=m\hbar$ (\tref{t:09}).
Этот пример показывает, что дискретность возникает из граничных условий.

\textbf{Гамильтониан} ($\op H\psi=E\psi$) — стационарное уравнение Шредингера (\tref{t:06}):
для связанных состояний спектр дискретен (яма, осциллятор, атом), для инфинитного
движения — непрерывен.

\subsection{Коммутаторы и одновременная измеримость}

Коммутатор $[\op A,\op B]=\op A\op B-\op B\op A$. Основное соотношение (вывод~\ref{der:5:comm})
\begin{equation}\label{eq:5:xp}
  [\op x,\op p_x]=\ii\hbar,\qquad [\op x,\op p_y]=0,\qquad [\op p_x,\op p_y]=0 .
\end{equation}
\textbf{Теорема.} Две величины имеют общую полную систему собственных функций (одновременно
измеримы, могут иметь одновременно точные значения) тогда и только тогда, когда их операторы
коммутируют. Если $[\op A,\op B]\neq0$, справедливо обобщённое соотношение неопределённостей
(вывод~\ref{der:5:robertson})
\begin{equation}\label{eq:5:rob}
  \Delta A\,\Delta B\ge\frac12\bigl|\avg{[\op A,\op B]}\bigr| ,
\end{equation}
которое для $x$ и $p_x$ даёт $\Delta x\,\Delta p_x\ge\hbar/2$ (\tref{t:02}).
Величина, оператор которой коммутирует с $\op H$ (и не зависит явно от времени), сохраняется:
её среднее не меняется со временем, а распределение вероятностей стационарно.

\section{Математические выводы}

\begin{derivation}[label=der:5:pherm]{эрмитовость оператора импульса}
Пусть $\varphi,\psi\to0$ при $x\to\pm\infty$ (нормируемые функции). Интегрируя по частям:
\[
  \int_{-\infty}^{\infty}\varphi^*\Bigl(-\ii\hbar\diff{\psi}{x}\Bigr)\dd x
  =-\ii\hbar\bigl[\varphi^*\psi\bigr]_{-\infty}^{\infty}+\ii\hbar\int\diff{\varphi^*}{x}\psi\,\dd x
  =\int\Bigl(-\ii\hbar\diff{\varphi}{x}\Bigr)^{\!*}\psi\,\dd x .
\]
Внеинтегральный член равен нулю, а $\ii\hbar\,\dd\varphi^*/\dd x=(-\ii\hbar\,\dd\varphi/\dd x)^*$.
Условие~\eqref{eq:5:herm} выполнено. Заметим: оператор $\dd/\dd x$ без множителя $\ii$ эрмитовым
не является (он антиэрмитов), поэтому множитель $-\ii$ в $\op p$ необходим. Эрмитовость $\op T$
следует из двукратного интегрирования по частям, $\op U$ — очевидна (умножение на
действительную функцию), следовательно, эрмитов и $\op H$.
\end{derivation}

\begin{derivation}[label=der:5:real]{собственные значения эрмитова оператора действительны}
Пусть $\op A\psi=a\psi$, $\int|\psi|^2\dd V=1$. Положим в~\eqref{eq:5:herm} $\varphi=\psi$:
\[
  \int\psi^*\op A\psi\,\dd V=a\int|\psi|^2\dd V=a,\qquad
  \int(\op A\psi)^*\psi\,\dd V=a^*\int|\psi|^2\dd V=a^* .
\]
Из эрмитовости левые части равны, значит $a=a^*$.
\end{derivation}

\begin{derivation}[label=der:5:orth]{ортогональность собственных функций}
Пусть $\op A\psi_1=a_1\psi_1$, $\op A\psi_2=a_2\psi_2$, $a_1\neq a_2$ (оба действительны).
Из~\eqref{eq:5:herm} с $\varphi=\psi_1$, $\psi=\psi_2$:
\[
  \int\psi_1^*\op A\psi_2\,\dd V=a_2\int\psi_1^*\psi_2\,\dd V,\qquad
  \int(\op A\psi_1)^*\psi_2\,\dd V=a_1\int\psi_1^*\psi_2\,\dd V .
\]
Вычитая: $(a_2-a_1)\int\psi_1^*\psi_2\,\dd V=0$, а так как $a_1\neq a_2$, то $\int\psi_1^*\psi_2\,\dd V=0$.
\end{derivation}

\begin{derivation}[label=der:5:avg]{среднее значение как статистическое среднее}
Подставим $\psi=\sum_nc_n\psi_n$ в~\eqref{eq:5:mean} (дискретный спектр, $\int\psi_m^*\psi_n=\delta_{mn}$):
\[
  \avg A=\int\Bigl(\sum_mc_m^*\psi_m^*\Bigr)\op A\Bigl(\sum_nc_n\psi_n\Bigr)\dd V
  =\sum_{m,n}c_m^*c_n\,a_n\int\psi_m^*\psi_n\,\dd V=\sum_n a_n|c_n|^2 .
\]
Это среднее значение случайной величины, принимающей значения $a_n$ с вероятностями $|c_n|^2$.
Нормировка $\psi$ даёт $\sum|c_n|^2=1$.
\end{derivation}

\begin{derivation}[label=der:5:pmean]{среднее импульса: обоснование $\op p=-\ii\hbar\,\dd/\dd x$}
По Борну, распределение импульса задаётся $|a(p)|^2$, где $\psi(x)=\int a(p)\frac{\ee^{\ii px/\hbar}}{\sqrt{2\pi\hbar}}\dd p$.
Подействуем на $\psi$ оператором $-\ii\hbar\,\dd/\dd x$: дифференцирование под интегралом даёт
множитель $p$:
\[
  -\ii\hbar\diff{\psi}{x}=\int p\,a(p)\frac{\ee^{\ii px/\hbar}}{\sqrt{2\pi\hbar}}\dd p .
\]
Тогда, используя $\int\ee^{\ii(p-p')x/\hbar}\dd x=2\pi\hbar\,\delta(p-p')$,
\[
  \int\psi^*\Bigl(-\ii\hbar\diff{\psi}{x}\Bigr)\dd x
  =\iint a^*(p')\,p\,a(p)\Bigl[\int\frac{\ee^{\ii(p-p')x/\hbar}}{2\pi\hbar}\dd x\Bigr]\dd p\,\dd p'
  =\int p\,|a(p)|^2\dd p=\avg p .
\]
Таким образом, формула~\eqref{eq:5:mean} с $\op p=-\ii\hbar\,\dd/\dd x$ даёт правильное среднее,
вычисляемое непосредственно в координатном представлении.
\end{derivation}

\begin{derivation}[label=der:5:comm]{коммутатор $[\op x,\op p_x]=\ii\hbar$}
Подействуем коммутатором на произвольную функцию $\psi(x)$:
\[
  [\op x,\op p_x]\psi=x\Bigl(-\ii\hbar\diff{\psi}{x}\Bigr)-\Bigl(-\ii\hbar\diff{}{x}\Bigr)(x\psi)
  =-\ii\hbar x\psi'+\ii\hbar(\psi+x\psi')=\ii\hbar\psi .
\]
Так как $\psi$ произвольна, $[\op x,\op p_x]=\ii\hbar$. Для $x$ и $p_y$: $\partial(x\psi)/\partial y=x\,\partial\psi/\partial y$,
поэтому $[\op x,\op p_y]=0$. Производные по разным переменным коммутируют, поэтому $[\op p_x,\op p_y]=0$.
\end{derivation}

\begin{derivation}[label=der:5:robertson]{соотношение неопределённостей $\Delta x\,\Delta p_x\ge\hbar/2$}
\emph{Допущение:} $\psi$ нормирована и убывает на бесконечности. Без ограничения общности
$\avg x=0$, $\avg{p_x}=0$ (иначе сделаем сдвиг $x\to x-\avg x$ и умножим $\psi$ на
$\ee^{-\ii\avg{p}x/\hbar}$; это не меняет $\Delta x$, $\Delta p$). Тогда $(\Delta x)^2=\avg{x^2}$,
$(\Delta p)^2=\avg{p^2}$.

Рассмотрим при действительном $\alpha$ неотрицательный интеграл
\[
  I(\alpha)=\int_{-\infty}^{\infty}\Bigl|\alpha x\psi+\hbar\diff{\psi}{x}\Bigr|^2\dd x\ge0 .
\]
Раскроем модуль:
\[
  I(\alpha)=\alpha^2\int x^2|\psi|^2\dd x+\alpha\hbar\int x\Bigl(\psi^*\diff{\psi}{x}+\psi\diff{\psi^*}{x}\Bigr)\dd x
  +\hbar^2\int\Bigl|\diff{\psi}{x}\Bigr|^2\dd x .
\]
\emph{Первый член:} $\alpha^2\avg{x^2}$.
\emph{Второй:} $\psi^*\psi'+\psi\psi^{*\prime}=\dd|\psi|^2/\dd x$; по частям
$\int x\,\dd|\psi|^2=[x|\psi|^2]-\int|\psi|^2\dd x=-1$. Итого $-\alpha\hbar$.
\emph{Третий:} по частям $\int|\psi'|^2\dd x=-\int\psi^*\psi''\dd x=\avg{p^2}/\hbar^2$, т.\,е.
третий член равен $\avg{p^2}$. Следовательно,
\[
  I(\alpha)=\avg{x^2}\alpha^2-\hbar\alpha+\avg{p^2}\ge0\quad\text{при всех }\alpha .
\]
Квадратный трёхчлен неотрицателен при всех $\alpha$, только если его дискриминант
неположителен: $\hbar^2-4\avg{x^2}\avg{p^2}\le0$, т.\,е.
\[
  \Delta x\,\Delta p_x\ge\frac{\hbar}{2}.
\]
Равенство достигается, если при некотором $\alpha$ $\alpha x\psi+\hbar\psi'=0$, т.\,е. $\psi\propto\ee^{-\alpha x^2/2\hbar}$ —
гауссов пакет. Для произвольных эрмитовых операторов аналогичное рассуждение даёт~\eqref{eq:5:rob}.
\end{derivation}

\section{Итоговые формулы}

\keyformula{\op{\vect r}=\vect r,\qquad \op{\vect p}=-\ii\hbar\nabla,\qquad \op T=-\frac{\hbar^2}{2m}\nabla^2,\qquad \op U=U(\vect r)}
 {$\nabla$ — оператор градиента, $\nabla^2$ — лапласиан.}
 {координатное представление, нерелятивистская частица без спина.}
 {построение операторов любых величин $A(\vect r,\vect p)$.}

\keyformula{\op H=\sum_i\Bigl(-\frac{\hbar^2}{2m_i}\nabla_i^2\Bigr)+U(\vect r_1,\dots,\vect r_N)}
 {$m_i$, $\nabla_i$ — масса и градиент по координатам $i$-й частицы; $U$ — полная потенциальная энергия.}
 {нерелятивистская система без спин-зависящих сил и магнитного поля.}
 {гамильтониан системы — входит в уравнение Шредингера.}

\keyformula{\op A\psi_n=a_n\psi_n,\qquad w_n=|c_n|^2,\ c_n=\int\psi_n^*\psi\,\dd V,\qquad \avg A=\int\psi^*\op A\psi\,\dd V}
 {$a_n$ — собственные значения, $\psi_n$ — ортонормированные собственные функции, $w_n$ — вероятность результата $a_n$.}
 {эрмитов оператор, дискретный спектр (для непрерывного — суммы заменяются интегралами).}
 {возможные результаты измерения и их вероятности; средние значения.}

\keyformula{[\op x,\op p_x]=\ii\hbar,\qquad \Delta A\,\Delta B\ge\tfrac12|\avg{[\op A,\op B]}|}
 {$[\op A,\op B]=\op A\op B-\op B\op A$, $\Delta$ — среднеквадратичные отклонения.}
 {любое состояние, для которого определены средние.}
 {критерий одновременной измеримости; соотношения неопределённостей.}

\section{Физическая интерпретация}

\textbf{Дискретность — следствие граничных условий.} Сами по себе дифференциальные уравнения на
собственные значения имеют решения при любых $a$; требование однозначности, конечности,
нормируемости отбирает дискретный набор (момент импульса, связанные состояния). Непрерывный
спектр соответствует инфинитному движению.

\textbf{Измерение.} Если состояние — собственная функция $\op A$, измерение $A$ даёт заведомо
одно значение. В суперпозиции результат случаен; повторное немедленное измерение даёт тот же
результат (состояние после измерения — собственная функция).

\textbf{Некоммутирующие величины.} Пары $x$ и $p_x$, $L_x$ и $L_y$ (\tref{t:09}) не могут иметь
одновременно точных значений. Коммутирующие наборы ($\op H,\op L^2,\op L_z$ для атома водорода)
дают квантовые числа, классифицирующие состояния (\tref{t:10}).

\section{Что необходимо знать на экзамене}

\begin{examcheck}
\checkhead{Определения}
\begin{itemize}
 \item линейный и эрмитов оператор; собственные функции и значения; спектр; вырождение;
 \item постулаты: результаты измерений — собственные значения; вероятности $|c_n|^2$; среднее $\int\psi^*\op A\psi$.
\end{itemize}
\checkhead{Формулы}
\begin{itemize}
 \item $\op x=x$, $\op p_x=-\ii\hbar\,\partial_x$, $\op T=-\hbar^2\nabla^2/2m$, $\op H=\op T+U$; гамильтониан системы;
 \item $[\op x,\op p_x]=\ii\hbar$; $\Delta A\Delta B\ge\frac12|\avg{[\op A,\op B]}|$.
\end{itemize}
\checkhead{Выводы}
\begin{itemize}
 \item эрмитовость $\op p$; действительность собственных значений; ортогональность собственных функций;
 \item $\avg A=\sum a_n|c_n|^2$; коммутатор $[x,p]$; доказательство $\Delta x\Delta p\ge\hbar/2$.
\end{itemize}
\checkhead{Интерпретации}
\begin{itemize}
 \item почему наблюдаемым соответствуют эрмитовы операторы; откуда берётся дискретность спектра;
 \item физический смысл коммутации; сохраняющиеся величины ($[\op A,\op H]=0$).
\end{itemize}
\end{examcheck}
